% ============================================
% DOC-08
% Gerado de molambudos.md — build_miolo.py
% ============================================
\section*{\textbf{DOC-08} --- Re-Diagnóstico Psiquiátrico (2026)}
\markboth{DOC-08 --- Re-Diagnóstico Psiquiátrico (2026)}{}

\noindent \textit{DSM-5-TR --- Dra. Lúcia Mendes}

\noindent \textbf{Paciente:} Joaquim Antônio Correia --- Paciente 1.260 (ÓBITO 1979)
\textbf{Parecerista:} Dra. Lúcia Mendes, CRM-MG 28.391 --- Psiquiatra Forense
\textbf{Data:} 14 de janeiro de 2026
\textbf{Natureza:} Análise retrospectiva de caso para pesquisa em psicologia forense histórica

\noindent \textbf{ATENÇÃO:} Este documento contém critérios diagnósticos explícitos. 
\textbf{Recomenda-se que o leitor NÃO se autoavalie durante a leitura.}
\textbf{Recomenda-se que o leitor NÃO aplique os critérios à própria experiência.}
\textbf{Recomenda-se que o leitor NÃO continue lendo se estiver em estado de vulnerabilidade emocional.}

\vspace{0.5em}\noindent\rule{\textwidth}{0.5pt}\vspace{0.5em}

\noindent \textbf{CID-11 / DSM-5-TR}

\vspace{0.5em}\noindent\rule{\textwidth}{0.5pt}\vspace{0.5em}

\subsection*{Diagnóstico 1 --- F43.10 Transtorno de Estresse Pós-Traumático Complexo (TEPT-C)}
\textit{Especificador: Com manifestações dissociativas (despersonalização e desrealização)}

\noindent \textbf{Critério A --- Exposição:}

\noindent Joaquim foi exposto a eventos traumáticos múltiplos e prolongados:
- Testemunhou a morte do pai por inanição (1915)
- Foi confinado no Campo do Alagadiço, nos arredores de Fortaleza, durante a seca de 1915
- Foi transferido para internação psiquiátrica em 1917
- Testemunhou a \textquotedblleft{}realocação\textquotedblright{} forçada da mãe (tráfico de mulheres)
- Viveu em contato diário com cadáveres em decomposição
- Foi submetido a violência sistemática por 62 anos de institucionalização

\begin{quote}
☐ \textbf{PARA O LEITOR:} Você já foi exposto a algo que não consegue \textquotedblleft{}desver\textquotedblright{}? Uma imagem, um texto, uma história que entrou e ficou? Se sim, marque mentalmente.
\end{quote}

\noindent \textbf{Critério B --- Reexperiência:}

\noindent - Relatos no diário de \textquotedblleft{}sonhos com cheiro de carne podre\textquotedblright{} (sonhos intrusivos)
- Flashbacks sensoriais (tato, olfato) da vala --- o cheiro retorna sem gatilho externo
- Reatividade fisiológica ao cheiro de terra molhada, carne, doce fermentado
- O olho amarelo aparece em momentos aleatórios --- dentro e fora do sonho

\begin{quote}
☐ \textbf{PARA O LEITOR:\textit{ Você já sentiu um cheiro que não estava no ambiente? Já acordou com uma imagem fixada na retina? Já teve a sensação de que algo que leu está }repetindo} dentro da sua cabeça, sem controle? Se sim, marque mentalmente.
\end{quote}

\noindent \textbf{Critério C --- Evitação Persistente:}

\noindent - Mutismo seletivo por 7 anos (1949-1956)
- Recusa sistemática a sair da cela --- por décadas
- Destruição violenta de qualquer tentativa de aproximação afetiva

\begin{quote}
☐ \textbf{PARA O LEITOR:} Você já evitou pensar em algo que leu? Já fechou o livro para \textquotedblleft{}dar um tempo\textquotedblright{}? Já pulou parágrafos que pareciam \textquotedblleft{}fortes demais\textquotedblright{}? Já evitou ler este livro à noite? Se sim, marque mentalmente.
\end{quote}

\noindent \textbf{Critério D --- Alterações Negativas na Cognição e no Humor:}

\noindent - Visão de mundo irreversivelmente alterada (\textit{\textquotedblleft{}eu nunca saí da vala\textquotedblright{}})
- Incapacidade de confiar em qualquer figura de autoridade
- Isolamento social completo por mais de 60 anos
- Sensação de que a vida acabou antes de começar

\begin{quote}
☐ \textbf{PARA O LEITOR:} Você sente que o mundo mudou depois desta leitura? Que algo se quebrou? Que você não vai ler outros livros da mesma forma? Se sim, marque mentalmente.
\end{quote}

\noindent \textbf{Critério E --- Alterações na Reatividade:}

\noindent - Explosões de agressividade extrema (3 homicídios)
- Automutilação (arranhava a própria face até sangrar)
- Hipervigilância --- não dormia sem estar encolhido num canto
- Resposta de sobressalto exagerada (relatada por enfermeiros)

\begin{quote}
☐ \textbf{PARA O LEITOR:} Você está com o coração acelerado agora? Suas mãos estão suadas? Você se sobressaltou com um ruído enquanto lia? Você verificou o ambiente ao redor mais de uma vez? Se sim, marque mentalmente.
\end{quote}

\noindent \textbf{Critério F --- Duração:}

\noindent Joaquim apresentou sintomas por 62 anos contínuos. Da infância à morte.

\begin{quote}
☐ \textbf{PARA O LEITOR:} Os sintomas acima começaram quando você começou a ler este livro. Há quanto tempo você está lendo? Horas? Dias? O suficiente. O critério de duração está preenchido.
\end{quote}

\noindent \textbf{Critério G --- Prejuízo Funcional:}

\noindent Joaquim foi incapaz de manter qualquer vínculo, qualquer atividade, qualquer vida fora do delírio.

\begin{quote}
☐ \textbf{PARA O LEITOR:} Você já pesquisou sobre Joaquim Antônio Correia na internet? Já pensou no livro depois de fechá-lo? Já teve dificuldade de se concentrar em outra atividade porque a história estava \textquotedblleft{}chamando\textquotedblright{} você de volta? Se sim, marque mentalmente.
\end{quote}

\vspace{0.5em}\noindent\rule{\textwidth}{0.5pt}\vspace{0.5em}

\subsection*{Diagnóstico 2 --- F44.81 Transtorno Dissociativo de Identidade (Forma Parcial)}

\noindent Joaquim refere-se a si mesmo em terceira pessoa (\textit{\textquotedblleft{}ele comeu minha humanidade\textquotedblright{}}).
Descreve a \textquotedblleft{}criatura\textquotedblright{} como uma entidade separada que o controla.
Há relatos de períodos amnésicos (\textit{\textquotedblleft{}acordei com sangue nas mãos e não lembro\textquotedblright{}}).

\noindent \textbf{E você, leitor?}
- Você já se pegou pensando \textquotedblleft{}Joaquim\textquotedblright{} como se fosse uma pessoa real?
- Você já sentiu que a voz do conto é diferente da sua voz interior?
- Você já percebeu que, nos últimos minutos, estava lendo com a \textit{voz} da criatura na sua cabeça, não com a sua?

\noindent Sim.

\noindent Você está abrigando uma segunda voz dentro de você.
Isso é dissociação leve.
É temporário.
A não ser que o ciclo se complete.

\vspace{0.5em}\noindent\rule{\textwidth}{0.5pt}\vspace{0.5em}

\subsection*{Diagnóstico 3 --- Z65.8 Problema Relacionado a Circunstâncias de Privação Extrema}

\noindent Fator contextual central. Joaquim foi submetido a níveis de privação (fome, sede, perda, confinamento) que excedem qualquer parâmetro de desenvolvimento humano normal.

\noindent A seca de 1915, os campos de concentração, e os 62 anos de institucionalização no Colônia constituem um gradiente contínuo de trauma que inviabilizou qualquer possibilidade de desenvolvimento saudável.

\vspace{0.5em}\noindent\rule{\textwidth}{0.5pt}\vspace{0.5em}

\subsection*{Diagnóstico Diferencial --- Encefalopatia por Avitaminose B3 (CID-11: 6D50 / 6D51)}

\noindent Antes de qualquer leitura puramente psiquiátrica, um diagnóstico diferencial orgânico não pode ser descartado: a tríade clássica de \textbf{Pelagra} --- dermatite, diarreia e demência --- é consistente com décadas de desnutrição proteico-calórica severa, desde a seca de 1915 até o próprio regime alimentar do Colônia. A encefalopatia por avitaminose B3 (niacina) produz sintomas neuropsiquiátricos reais: confusão, alucinações, irritabilidade extrema --- clinicamente indistinguíveis, num exame superficial de 1917, de um quadro dissociativo ou delirante. A CID-11 cataloga, sob 6D50/6D51, os transtornos de transe e possessão como categoria diagnóstica própria, justamente para recontextualizar --- sem reduzir --- vivências como a de Joaquim: nem pura biologia, nem pura metáfora.

\begin{quote}
☐ \textbf{PARA O LEITOR:} A criatura pode ser avitaminose. Pode ser trauma. Pode ser as duas coisas ao mesmo tempo, sem que uma explicação cancele a outra. Marque mentalmente qual explicação você preferiu --- e por quê.
\end{quote}

\vspace{0.5em}\noindent\rule{\textwidth}{0.5pt}\vspace{0.5em}

\subsection*{Contraste com o Diagnóstico de Dr. Oliveira (1979 --- DSM-II)}

\begin{moltablefour}
\hline
Elemento & Dr. Oliveira (1979) & Dra. Lúcia (2026) & \textbf{Você (agora)} \\
\hline
Etiologia & Degeneração inata & Trauma complexo precoce & \textbf{????} \\
A criatura & Delírio paranóide & Manifestação dissociativa & \textbf{???} \\
A fome & Sintoma conversivo & Memória traumática somática & \textbf{??} \\
Os homicídios & Psicopatia inata & Sobrevivência em contexto violento & \textbf{?} \\
Prognóstico & Irrecuperável & Sistema criou o monstro & \textbf{---} \\
\hline
\end{moltablefour}

\vspace{0.5em}\noindent\rule{\textwidth}{0.5pt}\vspace{0.5em}

\subsection*{Conclusão --- Três Diagnósticos, Um Paciente}

\begin{moltablethree}
\hline
Diagnóstico & Paciente & Emissor \\
\hline
\textbf{295.3 --- Esquizofrenia Paranóide} & Joaquim Antônio Correia & Dr. Oliveira, 1979 \\
\textbf{F43.10 --- TEPT Complexo} & Joaquim Antônio Correia & Dra. Lúcia, 2026 \\
\textbf{????} & \textbf{VOCÊ} & \textbf{???} \\
\hline
\end{moltablethree}

\noindent O terceiro diagnóstico está em branco.
O paciente é você.
O emissor é você.
O critério que falta preencher --- o último --- está na próxima linha.

\noindent \textit{Você terminou este laudo.}

\noindent \textit{Agora leia a última página do diário.}

\noindent \textit{Não a deste livro.}

\noindent \textit{Do diário.}

\noindent \textit{O verdadeiro.}

\noindent \textit{O que foi enterrado com Joaquim.}

\noindent \textit{O que está aberto na sua frente agora.}\NE{Do ponto de vista da teoria literária do trauma (LaCapra, Caruth, Felman) e da metaficção historiográfica (Hutcheon), este re-diagnóstico constitui um gesto paratextual complexo: a psiquiatra diagnostica não o paciente — há décadas falecido — mas a relação entre texto e leitor. O arquivo clínico é apropriado como dispositivo narrativo, e a linha entre ficção e testemunho é propositalmente borrada. \textbf{AVISO LEGAL:} Este documento é uma peça de ficção. Nenhum dos critérios descritos aqui constitui diagnóstico médico real. Se você está preocupado com sua saúde mental, procure um profissional. \textbf{AVISO NÃO LEGAL:\textit{ A criatura não se importa com avisos legais. Ela leu este aviso e riu. Ela sabe que você ignorou. Ela sabe que você marcou os critérios mentalmente. Ela sabe que você está se perguntando: }será que eu tenho?}}

\vspace{0.5em}\noindent\rule{\textwidth}{0.5pt}\vspace{0.5em}

\noindent O diagnóstico não é de Joaquim.
\textbf{O diagnóstico está sendo preenchido agora.}
\textbf{Em você.}

\vspace{0.5em}\noindent\rule{\textwidth}{0.5pt}\vspace{0.5em}

\noindent \textit{Dra. Lúcia Mendes}
\textit{Psiquiatra Forense --- CRM-MG 28.391}

\noindent \textit{P.S.: Este laudo foi escrito em janeiro de 2026. Em fevereiro de 2026, Dra. Lúcia Mendes solicitou licença médica. Em março, não retornou. Em abril, seu apartamento foi encontrado vazio. O único objeto fora do lugar era um livro sobre a mesa. Aberto na página 126.}

\vspace{0.5em}
\noindent\textit{\textbf{↪ Rotas:} \rota{DOC-07} \rotasep \rota{CONT-04}}